\section{Euler characteristic, signature, and finite covers}
\label{sec:consequences}

The concentration theorem determines the remaining $L^2$-Betti number
from an ordinary topological invariant. It also gives restrictions on
four-manifold signatures and, when sufficiently many finite covers
exist, on their rational homology.

\begin{corollary}\label{cor:euler}
Let $Q$ be a finite connected aspherical integral Poincar\'e complex
of formal dimension four, with arbitrary orientation character. Then
\[
 \chi(Q)=b_2^{(2)}(\widetilde Q)\ge0.
\]
Moreover, $\chi(Q)=0$ if and only if every $L^2$-Betti number of
$\widetilde Q$ vanishes. The same conclusions hold for closed connected
aspherical topological four-manifolds.
\end{corollary}
\begin{proof}
For a finite CW complex the $L^2$ Euler--Poincar\'e formula is
\[
 \chi(Q)=\sum_{p\ge0}(-1)^p b_p^{(2)}(\widetilde Q)
\]
\cite[Theorem~1.35(2)]{Luck2002}. Theorem~\ref{thm:main} leaves only
the degree-two term, whose von Neumann dimension is nonnegative.
The characterization of zero follows from the same theorem.
For a topological manifold, pass to the finite homotopy model of
Lemma~\ref{found:model}; both invariants are preserved.
\end{proof}

The middle term need not vanish: for closed oriented surfaces
$\Sigma_g,\Sigma_h$ of genera $g,h\ge2$, the aspherical product has
$\chi(\Sigma_g\times\Sigma_h)=4(g-1)(h-1)>0$.
The four-torus is an example with $\chi=0$.

For an oriented closed four-manifold $M$, let $\sigma(M)$ be the
signature of its real intersection pairing on $H_2(M;\mathbb R)$.
The next consequence is the Euler--signature inequality discussed by
Gromov and L\"uck; see \cite[Conjecture~1.3]{AlbaneseDiCerboLombardi2025}.

\begin{corollary}\label{cor:signature}
Every closed connected oriented aspherical topological four-manifold
satisfies
\[
 \chi(M)\ge |\sigma(M)|.
\]
In particular, $\chi(M)=0$ implies $\sigma(M)=0$.
\end{corollary}
\begin{proof}
The topological $L^2$-signature theorem gives
$\sigma(M)=\sigma^{(2)}(\widetilde M)$
\cite[Theorem~0.2]{LuckSchick2003}.
The latter is the difference of the positive and negative von Neumann
dimensions of the lifted middle-dimensional intersection form, so
$|\sigma^{(2)}(\widetilde M)|\le b_2^{(2)}(\widetilde M)$.
Corollary~\ref{cor:euler} now proves the inequality.
\end{proof}

This corollary includes nonsmoothable manifolds: the cited signature
theorem is topological. Its equality between ordinary and $L^2$
signatures is an additional input; it has not been assumed for
arbitrary finite Poincar\'e complexes.

\begin{corollary}\label{cor:betti-growth}
Let $Q$ be as in Corollary~\ref{cor:euler}, and suppose that
$\Gamma=\pi_1(Q)$ admits a descending sequence
$\Gamma=\Gamma_0\supseteq\Gamma_1\supseteq\cdots$ of finite-index
normal subgroups with $\bigcap_i\Gamma_i=\{1\}$.
Let $Q_i\to Q$ be the associated covers. For every $p\ge0$,
\[
 \lim_{i\to\infty}
 \frac{b_p(Q_i;\mathbb Q)}{[\Gamma:\Gamma_i]}
 =\begin{cases}
   \chi(Q),&p=2,\\
   0,&p\ne2.
  \end{cases}
\]
The same statement holds for a closed connected aspherical topological
four-manifold and its corresponding covers.
\end{corollary}
\begin{proof}
L\"uck's approximation theorem
\cite[Theorem~0.1]{Luck1994Approximation} identifies the limit with
$b_p^{(2)}(\widetilde Q)$. Apply Theorem~\ref{thm:main} and
Corollary~\ref{cor:euler}.
A finite homotopy model of a topological manifold lifts to a homotopy
equivalence on each corresponding finite cover, giving the last assertion.
\end{proof}
